\begin{abstract}
This paper adapts the Tsay-Lee framework to transactional memory models. Most modern concurrent algorithms are either lock-based or complex, hard-to-understand lock-free structures, while Software Transactional Memory (STM) remains underutilized. STM facilitates the non-blocking implementation of complex algorithms through transaction mechanisms. Originally designed to grant lock-free properties to Binary Search Trees, Tsay-Lee framework is used here to demonstrate how formal specifications can prove that universal algorithms maintain correctness when constructed from specialized components. The specifications are developed using the TLA+ language and the TLC model checker. The resulting specialization demonstrates that lock-based and non-blocking algorithms share more structural similarities than their source code suggests. This insight supports the creation of universal constructions for achieving lock-free properties. This work is intended for researchers seeking new universal paths to non-blocking algorithm implementations.
\end{abstract}
